\noindent\textbf{Elección del procedimiento.} En el núcleo $u_\theta=\Omega r$ y fuera $u_\theta=\Gamma/(2\pi r)$. La continuidad en $r=a$ fija $\Omega=\Gamma/(2\pi a^2)$. La vorticidad del giro rígido es $2\Omega$. La velocidad crece dentro y decrece fuera, por eso el máximo está en el borde.

\begin{enumerate}

\item Igualar las velocidades interior y exterior en el radio del núcleo.

\[\fitmath{\Omega=\frac{\Gamma}{2\cdot \pi\cdot {a}^{2}}}\]

\[\fitmath{\Omega=\frac{(400)}{2\cdot \pi\cdot {(3)}^{2}}\approx 7.073553\,\mathrm{1/s}}\]

\item La vorticidad es el doble de la velocidad angular del giro rígido.

\[\fitmath{\mathrm{wz}=2\cdot \Omega}\]

\[\fitmath{\mathrm{wz}=2\cdot (7.073553026)\approx 14.14711\,\mathrm{1/s}}\]

\item Evaluar la velocidad en el borde.

\[\fitmath{\mathrm{umax}=\Omega\cdot a}\]

\[\fitmath{\mathrm{umax}=(7.073553026)\cdot (3)\approx 21.22066\,\mathrm{m/s}}\]

\item Como 30 m supera 3 m, usar la rama exterior.

\[\fitmath{\mathrm{ur}=\frac{\Gamma}{2\cdot \pi\cdot r}}\]

\[\fitmath{\mathrm{ur}=\frac{(400)}{2\cdot \pi\cdot (30)}\approx 2.122066\,\mathrm{m/s}}\]

\end{enumerate}
